\documentclass[11pt]{article}

\usepackage[utf8]{inputenc}
\usepackage[T1]{fontenc}
\usepackage{amsmath, amssymb, bm}
\usepackage{geometry}
\geometry{margin=1.15in}

\title{Taylor-Proximal Approximation for Homoskedastic Matrix Noise}
\author{}
\date{}

\begin{document}
\maketitle

\section{Model}

We narrow the Taylor-proximal approximation to an isotropic Gaussian
likelihood with unknown variance.  Let
$\Omega \subseteq \{1,\ldots,m\}\times\{1,\ldots,n\}$ be the observed entries,
and let $g=\gamma^2$.  The model is
\begin{align}
    Y_{ij} \mid X_{ij}, g
        &\sim \mathcal{N}(X_{ij}, g),
        \qquad (i,j)\in\Omega, \\
    X \mid g,\lambda
        &\sim \mathrm{NND}\!\left(\frac{\lambda}{\sqrt{g}}\right), \\
    g &\sim P_g .
\end{align}
Equivalently, if $E=Y-X$, then $E_{ij}\mid g\sim\mathcal{N}(0,g)$ on
observed entries.  The scale factor $1/\sqrt{g}$ in the nuclear-norm prior is
the homoskedastic analogue of the scale coupling used by the Lasso: larger
noise weakens the effective penalty, while smaller noise strengthens it.  Since
the nuclear-norm distribution has normalizing constant proportional to
$\eta^{-mn}$ for parameter $\eta$, the prior density contributes
\begin{equation}
    \log p(X\mid g,\lambda)
        = \mathrm{const}
        + mn\log\!\left(\frac{\lambda}{\sqrt{g}}\right)
        - \frac{\lambda}{\sqrt{g}}\|X\|_* .
\end{equation}
Thus the negative log density contains both
$(\lambda/\sqrt{g})\|X\|_*$ and $(mn/2)\log g$.  This scaling is the piece that
keeps the posterior mode from degenerating toward the smallest admissible
variance.

For the implemented version we use an inverse-gamma variance prior,
\begin{equation}
    p(g) \propto g^{-(a_0+1)}\exp(-b_0/g),
\end{equation}
although the objective below separates the prior contribution so that another
one-dimensional prior on $g$ can be substituted.

\section{Taylor Delta Objective}

Let $\mu$ denote the variational posterior mean for $X$, and define
\begin{equation}
    B(\mu) = \left(\mu^\top \mu + \epsilon I\right)^{-1/2},
    \qquad
    D_i = \mathrm{diag}\!\left(\mathbf{1}\{(i,j)\in\Omega\}\right)_{j=1}^n .
\end{equation}
At fixed $(\mu,g)$ the row precision induced by the Taylor-delta approximation
is
\begin{equation}
    A_i(\mu,g)
        = \frac{1}{g}D_i
        + \frac{\lambda}{\sqrt{g}}B(\mu).
\end{equation}
Dropping constants independent of $(\mu,g)$, the collapsed minimization
objective is
\begin{align}
    \mathcal{J}(\mu,g)
        &=
        \frac{1}{2g}
        \sum_{(i,j)\in\Omega}(y_{ij}-\mu_{ij})^2
        + \frac{1}{2}\sum_{i=1}^m \log |A_i(\mu,g)| \notag\\
        &\quad
        + \frac{\lambda}{\sqrt{g}}\|\mu\|_*
        + \left(\frac{|\Omega|}{2} + \frac{mn}{2} + a_0 + 1\right)\log g
        + \frac{b_0}{g}.
        \label{eq:homoskedastic-objective}
\end{align}
The first $\log g$ term comes from the Gaussian likelihood, the second from the
normalizing constant of the scaled nuclear-norm prior, and the final two terms
from the inverse-gamma prior on $g$.

\section{Optimization}

The implementation optimizes $\mu$ and $g$ by alternating monotone steps.  For a
fixed variance $g$, split \eqref{eq:homoskedastic-objective} into a smooth part
$h_g(\mu)$ and the nonsmooth nuclear-norm term
$(\lambda/\sqrt{g})\|\mu\|_*$.  A proximal-gradient step is
\begin{align}
    Z^{(t)} &= \mu^{(t)} - s_t\nabla h_g(\mu^{(t)}), \\
    \mu^{(t+1)}
        &= \operatorname{SVT}_{s_t\lambda/\sqrt{g}}\!\left(Z^{(t)}\right),
\end{align}
where $\operatorname{SVT}_\tau$ soft-thresholds the singular values by $\tau$.
The step size $s_t$ is selected by backtracking until
$\mathcal{J}(\mu^{(t+1)},g)\leq\mathcal{J}(\mu^{(t)},g)$.

After the mean update, the scalar variance parameter is updated on the
unconstrained scale $\ell=\log g$:
\begin{equation}
    \ell^{(t+1)}
        = \ell^{(t)} - r_t
        \frac{\partial}{\partial \ell}
        \mathcal{J}\!\left(\mu^{(t+1)}, e^{\ell^{(t)}}\right),
\end{equation}
again with backtracking to enforce monotonic decrease of
$\mathcal{J}(\mu,e^\ell)$.  Optimizing $\log g$ makes positivity automatic and
keeps the variance update as a one-dimensional correction to the matrix
proximal step.

\section{Recovered Posterior Covariance}

At the final iterate, the row covariance approximation is recovered using the
same Taylor precision:
\begin{equation}
    \Sigma_i =
    \left[
        \frac{1}{\hat g}D_i
        + \frac{\lambda}{\sqrt{\hat g}}B(\hat\mu)
    \right]^{-1}.
\end{equation}
The returned effective penalty is
\begin{equation}
    \lambda_{\mathrm{eff}} = \frac{\lambda}{\sqrt{\hat g}},
\end{equation}
so downstream scoring code can reuse the same covariance and prior-variance
interfaces as the known-variance Taylor approximation.

\end{document}
